% ===== PRÉAMBULE CANONIQUE — à coller VERBATIM en tête de CHAQUE .tex =====
\documentclass[11pt,a4paper]{article}
\usepackage[utf8]{inputenc}\usepackage[T1]{fontenc}\usepackage{lmodern}
\usepackage[french]{babel}
\usepackage{amsmath,amssymb,mathtools}\usepackage{icomma}
\usepackage[version=4]{mhchem}          % \ce{...} : équations & formules
\usepackage{siunitx}\sisetup{locale=FR} % \qty{}{}, \si{}, \ang{}
\usepackage{chemfig}                    % \chemfig{...} : structures développées
\usepackage{xcolor}\usepackage{colortbl}
\usepackage{tikz}\usepackage{pgfplots}\pgfplotsset{compat=1.18}
\usepackage[most]{tcolorbox}\usepackage{enumitem}\usepackage{array,booktabs}
\usepackage[a4paper,margin=2.2cm]{geometry}
\usetikzlibrary{arrows.meta,calc,angles,quotes,positioning,patterns,backgrounds,shapes.geometric}
\definecolor{primary}{RGB}{222,49,99}\definecolor{primarydark}{RGB}{150,22,55}
\definecolor{defcol}{RGB}{0,114,178}\definecolor{methcol}{RGB}{52,92,148}
\definecolor{excol}{RGB}{0,135,90}\definecolor{attncol}{RGB}{200,40,55}
\tcbset{boxbase/.style={enhanced,breakable,boxrule=0.9pt,arc=2.5pt,
  left=3mm,right=3mm,top=2.3mm,bottom=2.3mm,fonttitle=\bfseries,coltitle=white}}
% --- boîtes TUTORIEL (noms EXACTS, ne pas renommer) ---
\newtcolorbox{cle}[1][]{boxbase,colback=defcol!6,colframe=defcol,
  title={\textsc{À retenir}\ifx\relax#1\relax\relax\else\textmd{ : \itshape #1}\fi}}
\newtcolorbox{methode}[1][]{boxbase,colback=methcol!7,colframe=methcol,sharp corners,
  title={\textsc{Méthode}\ifx\relax#1\relax\relax\else\textmd{ --- \itshape #1}\fi}}
\newtcolorbox{exemplebox}[1][]{boxbase,colback=excol!5,colframe=excol,
  title={\textsc{Exemple}\ifx\relax#1\relax\relax\else\textmd{ : \itshape #1}\fi}}
\newtcolorbox{piege}[1][]{boxbase,colback=attncol!6,colframe=attncol,
  title={\textsc{Piège}\ifx\relax#1\relax\relax\else\textmd{ : \itshape #1}\fi}}
\newtcolorbox{remarque}[1][]{boxbase,colback=black!4,colframe=black!45,
  title={\textsc{Remarque}\ifx\relax#1\relax\relax\else\textmd{ : \itshape #1}\fi}}
% --- boîtes EXERCICE / CORRIGÉ (noms EXACTS, auto-comptées) ---
\newtcolorbox[auto counter]{exo}[1][]{boxbase,colback=primary!4,colframe=primary,
  title={\textsc{Exercice}~\thetcbcounter\ifx\relax#1\relax\relax\else\textmd{ --- \itshape #1}\fi}}
\newtcolorbox[auto counter]{corrige}[1][]{boxbase,colback=excol!4,colframe=excol,
  title={\textsc{Corrigé}~\thetcbcounter\ifx\relax#1\relax\relax\else\textmd{ --- \itshape #1}\fi}}
\newcommand{\R}{\mathbb{R}}\newcommand{\N}{\mathbb{N}}\newcommand{\Z}{\mathbb{Z}}\newcommand{\ds}{\displaystyle}
% (un \usepackage{fancyhdr}+en-tête est possible ; il est ignoré côté web)
% =========================================================================
\begin{document}

% THEME_TITLE: Masse volumique, densité et volume molaire des gaz
\noindent
On ne pèse pas toujours : souvent on \emph{mesure un volume}. Ce thème installe les ponts entre
\textbf{masse} et \textbf{volume}, pour les liquides (masse volumique, densité) comme pour les gaz
(volume molaire).

\begin{cle}[masse volumique $\rho$]
La \textbf{masse volumique} $\rho$ est la masse d'une unité de volume :
\[ \boxed{\,\rho = \dfrac{m}{V}\,}\qquad\Longleftrightarrow\qquad m = \rho\,V. \]
Unités usuelles : \si{\gram\per\milli\litre} $=$ \si{\gram\per\cubic\centi\metre}. \textbf{Astuce} :
« $\rho = x\ \si{\gram\per\milli\litre}$ » signifie que $\SI{1}{\milli\litre}$ pèse $x\ \si{\gram}$,
donc que $\SI{1}{\litre}$ pèse $1000\,x\ \si{\gram}$. On part presque toujours « d'un litre ».
\end{cle}

\begin{cle}[densité $d$]
La \textbf{densité} $d$ est le rapport de la masse volumique du corps à celle de l'eau
($\rho_{\text{eau}} = \SI{1}{\gram\per\milli\litre}$) :
\[ d = \dfrac{\rho_{\text{corps}}}{\rho_{\text{eau}}}\quad(\text{sans unité}). \]
Comme $\rho_{\text{eau}} = \SI{1}{\gram\per\milli\litre}$, la densité a la \emph{même valeur} que $\rho$
exprimée en \si{\gram\per\milli\litre} : $d = 0{,}8 \Leftrightarrow \rho = \SI{0.8}{\gram\per\milli\litre}$.
\end{cle}

\begin{methode}[le cas des gaz : le volume molaire $V_m$]
À température et pression fixées, \emph{une mole de n'importe quel gaz parfait occupe le même volume}
$V_m$ :
\[ \boxed{\,V = n\times V_m\,}\qquad\Longleftrightarrow\qquad n = \dfrac{V}{V_m}. \]
Aux \textbf{CNTP} ($\SI{0}{\celsius}$, $1$~atm) : $V_m = \SI{22.4}{\litre\per\mole}$ ; dans les
conditions standard ($\SI{25}{\celsius}$) : $V_m \approx \SI{24.4}{\litre\per\mole}$.
\end{methode}

\begin{center}
\begin{tikzpicture}[scale=1,
   box/.style={draw,rounded corners=3pt,align=center,font=\footnotesize,minimum height=1cm,minimum width=2.7cm,line width=0.9pt}]
  \node[box,draw=defcol,fill=defcol!8] (V) at (0,0) {Volume gaz\\ $V$ (\si{\litre})};
  \node[box,draw=excol,fill=excol!8] (n) at (4.7,0) {Quantité\\ $n$ (\si{\mole})};
  \node[box,draw=primary,fill=primary!8] (m) at (9.4,0) {Masse\\ $m$ (\si{\gram})};
  \draw[-{Stealth[length=2mm]},line width=0.9pt] ([yshift=4pt]V.east) -- node[above,font=\scriptsize]{$\div V_m$} ([yshift=4pt]n.west);
  \draw[-{Stealth[length=2mm]},line width=0.9pt] ([yshift=-6pt]n.west) -- node[below,font=\scriptsize]{$\times V_m$} ([yshift=-6pt]V.east);
  \draw[-{Stealth[length=2mm]},line width=0.9pt] ([yshift=4pt]n.east) -- node[above,font=\scriptsize]{$\times M$} ([yshift=4pt]m.west);
  \draw[-{Stealth[length=2mm]},line width=0.9pt] ([yshift=-6pt]m.west) -- node[below,font=\scriptsize]{$\div M$} ([yshift=-6pt]n.east);
\end{tikzpicture}
\end{center}

\begin{exemplebox}[du volume à la masse d'un gaz]
Masse de $\SI{5.6}{\litre}$ de \ce{CO2} aux CNTP ($M=\SI{44}{\gram\per\mole}$) :
$n = \dfrac{5{,}6}{22{,}4} = \SI{0.25}{\mole}$, puis $m = 0{,}25\times 44 = \SI{11}{\gram}$.
\end{exemplebox}

\begin{piege}[deux confusions à éviter]
\begin{itemize}[leftmargin=1.4em,topsep=1pt,itemsep=1pt]
  \item $\SI{22.4}{\litre\per\mole}$ n'est valable qu'\textbf{aux CNTP} ($\SI{0}{\celsius}$) :
        à $\SI{25}{\celsius}$ une mole occupe $\approx\SI{24.4}{\litre}$.
  \item $\rho$ (masse \emph{totale} de solution par volume) $\neq$ concentration massique $C_m$
        (masse du seul \emph{soluté} par volume).
\end{itemize}
\end{piege}

\end{document}
